\documentclass[10pt]{mypackage}

\usepackage[uprightgreeks,light]{kpfonts}
\renewcommand{\coloneq}{\coloneqq}
%\renewcommand{\mathbb}{\mathds}
\usepackage{homework}
%\usepackage{notes}

%\usepackage[ backend=bibtex, style = alphabetic, sorting=ynt ]{biblatex}
%\addbibresource{  }

\usepackage{parskip}

\fancyhf{}
\fancyhead[R]{Avinash Iyer}
\fancyhead[L]{Algebraic Topology: Homework 28}
\fancyfoot[C]{\thepage}

\setcounter{secnumdepth}{0}

\begin{document}
\RaggedRight
\section{Revised Problems}%
\begin{problem}[Homework 21, Problem 1 (a)]
  Prove that $H_0\left( X,A \right) = 0$ if and only if $A$ intersects each path component of $X$ nontrivially.
\end{problem}
\begin{solution}
  Suppose $H_0\left( X,A \right) = 0$. From the long exact sequence, the map $i_{\ast}\colon H_0\left( A \right)\rightarrow H_0\left( X \right)$ is surjective, so the homology classes $\left[ \alpha \right]\in H_0\left( A \right)$ corresponding to path components of $A$ generate $H_0\left( X \right)$ upon inclusion into $X$. Therefore, the collection of path components of $A$ surjects onto the path components of $X$, so the intersection of $A$ with every path component of $X$ is nontrivial.

  In the reverse direction, if $A$ intersects every path component of $X$ nontrivially, then there are homology classes $\left[ \alpha \right]\in H_0\left( A \right)$ that correspond to path components of $H_0\left( X \right)$ via inclusion of $A$ into $X$, meaning $i_{\ast}\colon H_0\left( A \right)\rightarrow H_0\left( X \right)$ is surjective. Therefore, the map $j_{\ast}\colon H_0\left( X \right)\rightarrow H_0\left( X,A \right)$ has $\ker\left( j_{\ast} \right) = \img\left( i_{\ast} \right)$, so $j_{\ast}$ is the zero map. Since $\partial\colon H_0\left( X,A \right)\rightarrow H_{-1}\left( A \right) = 0$ has $\ker\left( \partial \right) = H_{0}\left( X,A \right)$, and $\ker\left( \partial \right) = \img\left( j_{\ast} \right)$, it follows that $H_0\left( X,A \right) = 0$.
\end{solution}
\begin{problem}[Homework 26, Problem 1]
  Prove that $d_q = j_{\ast}\circ \partial$.
\end{problem}
\begin{solution}
  We consider the following diagram defined by the exact sequences of the pair $\left( X^q,X^{q-1} \right)$ along the diagonal and $\left( X^{q-1},X^{q-2} \right)$ along the vertical.
 \begin{center}
   % https://tikzcd.yichuanshen.de/#N4Igdg9gJgpgziAXAbVABwnAlgFyxMJZARgBoBmAXVJADcBDAGwFcYkQAdDgYyghwQBfUuky58hFACYK1Ok1bsAEgH1gAR0EAKABoA9daX0aAtMUEBKEMNHY8BIuVk0GLNohCrT53Xu-DjdRMpS2sREAw7CSIAFmd5N3YuXn4hcMjxBxQABlIAVjlXRQ9ssNtMyRJSGMKFd08VdV91KxsIsXtKpylaxI8vIJ9As1C2jM7Y0mJe4ob-Ztb0juiUPNJsmfrSseWs5Cdplzr2baWovacao765oJDfUxDF8omUJwLr2e25GCgAc3gRFAADMAE4QAC2SFyIBwECQ5nCYMhCJocKQMgSsygjTKIGRUMQa1h8MQZCx9QAVmouPQ4DhBHiCUg4iSkHk2syyWjSQA2T71LhoeigvBMJngwkATh5SF5nMlctliExRSpNI4dIZEpRKuVxGyCt1Bv1iJBisQ-LZiAA7EbCZj0YhyPakDblQAOQSUQRAA
\begin{tikzcd}
             &                                 &                                                              &                                                   &                         & 0 \\
             &                                 &                                                              & 0 \arrow[d]                                       & H_{q-1}(X^q) \arrow[ru] &   \\
             &                                 &                                                              & H_{q-1}(X^{q-1}) \arrow[d, "j_{\ast}"] \arrow[ru] &                         &   \\
             & \cdots \arrow[r]                & {H_{q}(X^q,X^{q-1})} \arrow[r, "d_q"] \arrow[ru, "\partial"] & {H_{q-1}(X^{q-1},X^{q-2})} \arrow[d] \arrow[r]    & \cdots                  &   \\
             & H_q(X^q) \arrow[ru, "j_{\ast}"] &                                                              & H_{q-2}(X^{q-2}) \arrow[d]                        &                         &   \\
0 \arrow[ru] &                                 &                                                              & 0                                                 &                         &  
\end{tikzcd}
 \end{center} 
 This allows us to consider all instances of the map $d_q$ as arising from these respective exact sequences; i.e., the map $d_{q-1}$ arises from the exact sequences of the pairs $\left( X^{q-1},X^{q-2} \right)$  and $\left( X^{q-2},X^{q-3} \right)$ fitting into the above diagram, meaning that $d_q = j_{\ast}\circ\partial$.
\end{solution}
\section{Current Problems}%
\begin{problem}[Problem 1]
  Recall from class that for a CW complex $X$, the map
  \begin{align*}
    d_{q+1}\colon C_{q+1}^{\operatorname{CW}}\left( X \right)\rightarrow C_{q}^{\operatorname{CW}}\left( X \right)
  \end{align*}
  is given by
  \begin{align*}
    d_{q+1}\left( e_{i}^{q+1} \right)&= \sum_{j=1}^{\ell}n_{ij}e_j^{q}
  \end{align*}
  for some $n_{ij}\in \Z$, where $e_{i}^{q+1}$ is a $\left( q+1 \right)$-cell of $X$ and $e_j^{q}$ is a $q$-cell. Prove that each coefficient $n_{ij}$ is the degree of the composite map
  \begin{align*}
    S^{q}\xrightarrow{\varphi_i}X^q\rightarrow X^q/X^{q-1}\cong \bigvee_{j=1}^{\ell}S_j^{q}\rightarrow S_j^{q}.
  \end{align*}
\end{problem}
\begin{solution}
  We let $\Phi_i$ be the characteristic map for $D^{q+1}$ corresponding to $\varphi_i\colon S^q\rightarrow X^q$. This maps $D^{q+1}$ to $e^{q+1}_{j}$. By the long exact sequence of the pair $\left( D^{q+1},S^q \right)$, we observe that
  \begin{align*}
    H_{q+1}\left( S^q \right)\rightarrow H_{q+1}\left( D^{q+1},S^q \right)\rightarrow H_{q}\left( S^q \right)\rightarrow H_q\left( D^{q+1} \right),
  \end{align*}
  which induces an isomorphism $ H_{q+1}\left( D^{q+1},S^q \right)\cong H_q\left( S^q \right) $. Composing these maps with $d_{q+1}$ and the various quotient maps, we have the following diagram.
  \begin{center}
    
% https://tikzcd.yichuanshen.de/#N4Igdg9gJgpgziAXAbVABwnAlgFyxMJZARgBoAGAXVJADcBDAGwFcYkQAJAfQEcAKAMoA9HgEoQAX1LpMufIRQAmCtTpNW7bsB4BqYhL4ARIdr1ThYydJAZseAkQDMKmgxZtEnLqf18AGia6+qQBllIydvJOpMSqbhqeWjwGoSGBALT64uE2svYKyM6KceoeIAA65QDuWLB4jLDAHBLeyf4iAPQB2pkS2da2cg4o5C5q7uyVNXVYDTBNLdoGFv0RQwXKxa6lk+UARlgA5hBoLHDeAFYAvPomlTCMjBJTtTD1jc2tyyKruZHDJFIW3GCQq1Ve73mnyWgkCiwufUkqhgUEO8CIoAAZgAnCAAWyQoxAOAgSGCWNxBMQZGJpMQyhBZT4lQACgALLBcLCibyVehwHASKwU-FIBkkpDORnsKCtMzCkA40WIKUSxAAFm2E08aF55X5goVSqpAFYaGryDljUhNbSkAA2K2Uh3mukAdi1oOZ5TQnIuPOAfIFQqdyrNdsQHulnkqAGMoBAcAACMDeLAIpESIA
\begin{tikzcd}
\widetilde{H}_{q}\left(S^q\right) \arrow[r] \arrow[rdd, "\cdot n_{ij}"] & H_q\left(S^q\right) \arrow[r]           & {H_{q+1}\left(D^{q+1},S^q\right)} \arrow[r, "\left(\Phi_i\right)_{\ast}"]                      & {H_{q+1}\left(X^{q+1},X^q\right)} \arrow[d, "d_{q+1}"] \\
                                                             &                              &                                                                          & {H_{q}\left(X^q,X^{q-1}\right)} \arrow[d, "p_{\ast}"]  \\
                                                             & \widetilde{H}_{q}\left(S^{q}_{j}\right) & \bigoplus_{j=1}^{\ell}\widetilde{H}_{q}\left(S^q\right) \arrow[l, "\left(\pi_j\right)_{\ast}"] & \widetilde{H}_{q}\left(X^q/X^{q-1}\right) \arrow[l]   
\end{tikzcd}
  \end{center}
  This map extracts the coefficient on $e_{j}^{q}$ as $\left( \pi_j \right)_{\ast}$ is the projection onto $e_j^{q}$ in the direct sum decomposition of $ \widetilde{H}_q\left( X^q/X^{q-1} \right) $. This necessarily yields the degree of the map from $ \widetilde{H}_{q}\left( S^q \right)\rightarrow \widetilde{H}_q\left( S^q_{j} \right) $.
\end{solution}
\begin{problem}[Problem 2]
  Suppose that the following diagram of abelian groups is commutative with exact rows and the indicated maps isomorphisms.
  \begin{center}
    % https://tikzcd.yichuanshen.de/#N4Igdg9gJgpgziAXAbVABwnAlgFyxMJZABgBpiBdUkANwEMAbAVxiRAB12BjKCHBAL6l0mXPkIoAjOSq1GLNgGEA+sACOAakkCQQkdjwEiAJhnV6zVohABBZWt3CQGA+KIBmM3MtsAQvcd9MSMUABYvCwVrFQc9Z1FDCWQAVgj5K1tVNQBabUD41xDkADY0n2tOHj5BJxdgpLJJWUiMyt5+fLrEomkm83SlLK0BAHJOhLcUUz7vKMz1UfHCpM8Zlr8sxbiuyeRwtYHo+zHtiaLUg-L5nO0T2rOk0su5turdWRgoAHN4IlAAMwAThAALZIaQgHAQJAADn6V0qBC++SBoKQpkh0MQAE44qiwYgIVD0fCXuw0HRAXhGCjgQSMcTEJ5Zhl-rS0UzqIzwiy2Mi8XSkDzGaleRVyZTqQx2QS4ZikNjSa0JVSsIw7gDBTiuVjJGQxSB-hrDVq9TrwRD1tYvsb8UhmYy9QKOcLddonHbCUTdcZnQTRY7fR6tfrHX6kAB2c2IGHhwkM3XuOMBrHFAQUARAA
    \begin{tikzcd}
      \cdots \arrow[r] & C_{q+1} \arrow[d, "\cong"] \arrow[r, "\partial"] & A_q \arrow[d, "a_q"] \arrow[r, "f_q"] & B_q \arrow[r, "g_q"] \arrow[d, "b_q"] & C_q \arrow[r, "\partial"] \arrow[d, "\cong"] & A_{q-1} \arrow[d, "a_{q-1}"] \arrow[r] & \cdots \\
    \cdots \arrow[r] & C_{q+1}' \arrow[r, "\partial'"]                  & A_{q}' \arrow[r, "f_q'"]       & B_{q}' \arrow[r, "g_q'"]       & C_q' \arrow[r, "\partial'"]                      & A_{q-1}' \arrow[r]          & \cdots
    \end{tikzcd}
  \end{center}
  Show that there is a long exact sequence of the following form.
  \begin{center}
    % https://tikzcd.yichuanshen.de/#N4Igdg9gJgpgziAXAbVABwnAlgFyxMJZABgBpiBdUkANwEMAbAVxiRAB12BjKCHBAL6l0mXPkIoAjOSq1GLNgEEA+gEcQQkdjwEiAJhnV6zVohArgqgQHJOENMzgACAEJqNwkBm3iiAZkM5EzY3VWsPLTFdFAAWQOMFMwtVAFpJAQivUR0JZABWePlTc2VLNJs7ByZnNzL0zO8o3IA2QuCzWtT08M0sn2jkAHY2xI5uXn4NWRgoAHN4IlAAMwAnCABbJDIQHAgkdM9Vjf3qXaQ9XqPNxAMdvcQ-S7XrgLukGKfjxDi3xDzP64FX7NAFIVq-QYCCgCIA
    \begin{tikzcd}
    \cdots \arrow[r] & A_q \arrow[r] & A_{q}'\oplus B_q \arrow[r] & B_q' \arrow[r] & A_{q-1} \arrow[r] & A_{q-1}'\oplus B_{q-1} \arrow[r] & B_{q-1}' \arrow[r] & \cdots
    \end{tikzcd}
  \end{center}
\end{problem}
\begin{solution}
  Label the maps as $a_q\colon A_q\rightarrow A_q'$, $b_q\colon B_q\rightarrow B_q'$, and $c_q\colon C_q\rightarrow C_{q}'$.

  We claim that the maps given by
  \begin{align*}
    r_q &= a_q\oplus f_q\\
    s_q &= f_q'\circ \pi_1 - b_q\circ \pi_2,\\
    t_q &= \partial\circ c_q^{-1}\circ g_q'
  \end{align*}
  where $\pi_i$ denotes the projection onto coordinate $i$, yield the long exact sequence
  \begin{align*}
    \cdots\rightarrow A_q \xrightarrow{r_q} A_q'\oplus B_q \xrightarrow{s_q} B_q' \xrightarrow{t_q} A_{q-1}\rightarrow\cdots
  \end{align*}
  This requires us to show exactness at every level.

  First, let $a\in A_q$. Then we have $r_q(a) = \left( a_q(a),f_q(a) \right)$, so we get
  \begin{align*}
    s_q\circ r_q(a) &= f_q'\circ a_q(a) - b_q\circ f_q(a).
  \end{align*}
  By commutativity, we have $f_q'\circ a_q = b_q\circ f_q$, so thus $\img\left( r_q \right)\subseteq \ker\left( s_q \right)$. Now, suppose $\left( a',b \right)\in \ker\left( s_q \right)$, so that
  \begin{align*}
    0 &= f_q'\left( a' \right) - b_q(b).
  \end{align*}
  Applying $g_q'$, we have
  \begin{align*}
    0 &= g_q'\circ f_q'\left( a' \right) - g_q'\circ b_q(b),
  \end{align*}
  whence $g_q'\circ b_q(b) = 0$ by the exactness of the bottom row. Applying commutativity gives that $c_q\circ g_q(b) = 0$, meaning that $g_q(b) = 0$ since $c_q$ is an isomorphism, so $b\in \ker\left( g_q \right)$, so there is $a$ such that $f_q(a) = b$. We claim that $a_q(a) - a'\in \ker\left( f_q \right)$. For this, we take
  \begin{align*}
    f_q'\left( a_q(a) - a' \right) &= f_q'\circ a_q(a) - f_q'\left( a' \right)\\
                                   &= b_q\left( f_q(a) \right) - f_q'\left( a' \right)\\
                                   &= b_q\left( b \right) - f_q'\left( a' \right)\\
                                   &= 0,
  \end{align*}
  where we use commutativity in the second equality. By exactness, there is some $c'\in C_{q+1}'$ such that $\partial'\left( c' \right) = a_q(a) - a$. As $c_{q+1}$ is an isomorphism we have unique $c\in C_{q+1}$ $c_{q+1}\left( c \right) = c'$. It then follows that
  \begin{align*}
    a_q(a) - a' &= \partial'\left( c' \right)\\
                &= \partial'\circ c_{q+1}\left( c \right)\\
                &= a_{q}\circ \partial\left( c \right),
  \end{align*}
  so we have $a_q\left( a-\partial(c) \right) = a'$. Since $f_q\left( a-\partial(c) \right) = f_q(a)$ by exactness, we have that $r_q\left( a-\partial(c) \right) = \left( a',b \right)$, whence $\ker\left( s_q \right) = \img\left( r_q \right)$.

  Now, let $\left( a',b \right)\in A_q'\oplus B_q$. Then, we have $s_q\left( a',b \right) = f_q'\left( a' \right)-b_q(b)$. Therefore,
  \begin{align*}
    t_q\circ s_q\left( a',b \right) &= \partial\circ c_q^{-1}\circ g_q'\circ f_q'\left( a \right) - \partial\circ c_q^{-1}\circ g_q'\circ b_q(b)\\
                                    &= \partial\circ c_q^{-1}\circ c_q\circ g_q(b)\\
                                    &= \partial\circ g_q(b)\\
                                    &= 0,
  \end{align*}
  where we use exactness and commutativity to reduce the compositions, so that $\img\left( s_q \right)\subseteq \ker\left( t_q \right)$.

  Now, suppose $b'\in \ker\left( t_q \right)$. This gives
  \begin{align*}
    0 &= \partial\circ c_{q}^{-1}\circ g_q'\left(b'\right).
  \end{align*}
  If we let $c = c_q^{-1}\circ g_q'\left( b' \right)$, then we observe that $\partial(c) = 0$, so there is some $b\in B_q$ such that $g_q(b) = c$. We have
  \begin{align*}
    g_q'\circ b_q(b) &= c_q\circ g_q(b)\\
                     &= c_q\circ c_q^{-1}\circ g_q'\left( b' \right)\\
                     &= g_q'\left( b' \right),
  \end{align*}
  so that $b_q(b) - b'\in \ker\left( g_q' \right)$. There is then some $a'\in A_q'$ such that $f_q'\left( a' \right) = b_q(b)-b'$, or that $f_q'\left( a' \right)-b_q(b) = -b'$. We observe then that
  \begin{align*}
    s_q\left( -a',-b \right) &= f_q'\left( -a' \right) - b_q\left( -b \right)\\
                             &= b',
  \end{align*}
  so that $\ker\left( t_q \right)= \img\left( s_q \right) $.

  Finally, let $b'\in B_q'$. Then, we have
  \begin{align*}
    r_{q-1}\circ t_q\left( b' \right) &= r_{q-1}\left( \partial\circ c_q^{-1}\circ g_q'\left( b' \right) \right)\\
                                      &= \left( a_{q-1}\circ \partial \circ c_q^{-1}\circ g_q'\left( b' \right),f_{q-1}\circ \partial \circ c_q^{-1}\circ g_q'\left( b' \right) \right)\\
                                      &= \left( \partial'\circ c_q\circ c_q^{-1}\circ g_q'\left( b' \right),0 \right)\\
                                      &= \left( 0,0 \right),
  \end{align*}
  so that $\img\left( t_q \right)\subseteq \ker\left( r_{q-1} \right)$.

  Suppose $a\in \ker\left( r_{q-1} \right)$. Then, we have 
  \begin{align*}
    r_{q-1}\left( a \right) &= \left( a_{q-1}(a),f_{q-1}(a) \right)\\
                            &= 0.
  \end{align*}
  By exactness, there is some $c\in C_q$ such that $\partial (c) = a$. We have
  \begin{align*}
    a_{q-1}\circ \partial(c) &= \partial'\circ c_q\left( c \right)\\
                             &= 0,
  \end{align*}
  meaning that $c_q\left( c \right) = c'\in \ker\left( \partial' \right)$, so there is some $b'\in B_q$ such that $g_q'\left( b' \right) = c'$. We thus get
  \begin{align*}
    t_q\left( b' \right) &= \partial\circ c_q^{-1}\circ g_q'\left( b' \right)\\
                         &= \partial\circ c_q^{-1}\left( c' \right)\\
                         &= \partial\circ c_q^{-1}\circ c_q\left( c \right)\\
                         &= \partial(c)\\
                         &= a.
  \end{align*}
  Thus, $\ker\left( r_{q-1} \right)\subseteq \img\left( t_q \right)$, and we are done.
\end{solution}
\end{document}
